% ============================================================================
% Mục 4.4 - Hiện thực mã nguồn của phần B (features.py) và C (evaluate.py)
% ============================================================================
\subsection{Hiện thực mã nguồn (Phần B và C)}\label{sec:impl}

Mục này mô tả chi tiết hai mô-đun dùng chung do phần B và C viết: \code{src/features.py} (đặc trưng) và
\code{src/evaluate.py} (đánh giá). Với mỗi file, báo cáo liệt kê các thư viện được import, các biến toàn cục
(global), các lớp và hàm, kèm giải thích đầu vào, đầu ra và cách hoạt động. Mô-đun huấn luyện
\code{train_models.py} chỉ \emph{gọi} hai mô-đun này nên không được mô tả lại ở đây.

\subsubsection{Tổng quan luồng gọi hàm}\label{sec:impl-flow}

\begin{lstlisting}[language={},numbers=none,caption={Luồng gọi hàm từ dữ liệu đã chia đến bảng kết quả.},label={lst:callgraph}]
features.make_datasets(cfg)
 |-- load_splits()                    -> {train, val, test} (Parquet from part A)
 |-- build_feature_frame(full, cfg)   -> time + target + TSO + 39 features
 |    |-- add_calendar_features(full)       15 columns cal_*
 |    |-- add_lag_features(full, cfg)       13 columns lag_*, roll_*, diff_24
 |    `-- add_weather_features(full, cfg)   11 columns w_*
 |-- cut back by train/val/test time bounds, dropna (lag warm-up)
 `-- feature_columns(frame)           -> list of model inputs

train_models.py: fit model on feature_columns -> y_pred
evaluate.ResultsLog.add(model, split, y_true, y_pred)
 `-- evaluate() -> mae, rmse, mape, smape, r2, max_ae  (each calls _as_arrays)
evaluate.plot_* / error_by_group  -> figures in plots/
\end{lstlisting}

\subsubsection{Các thư viện được import}\label{sec:impl-imports}

Hai mô-đun \textbf{không thêm thư viện bên thứ ba nào mới}: mọi gói đều đã có trong \code{pyproject.toml}
của phần A, nên \code{uv sync} không cần thay đổi. Các import ``mới'' (chưa xuất hiện trong
\code{data_preprocessing.py}) là \code{from __future__ import annotations} và \code{dataclasses.field}.

\begingroup\footnotesize
\begin{longtable}{>{\raggedright\arraybackslash}p{0.30\linewidth}>{\raggedright\arraybackslash}p{0.13\linewidth}>{\raggedright\arraybackslash}p{0.07\linewidth}p{0.40\linewidth}}
\caption{Các import trong \code{features.py} (F) và \code{evaluate.py} (E).}\label{tab:imports}\\
\toprule
Import & Dùng ở & Mới? & Vai trò\\
\midrule
\endfirsthead
\toprule
Import & Dùng ở & Mới? & Vai trò\\
\midrule
\endhead
\code{from __future__ import annotations} & F, E & Có &
Hoãn việc đánh giá chú thích kiểu (type hint): Python lưu chúng dưới dạng chuỗi thay vì thực thi ngay. Nhờ đó
có thể viết gọn \code{FeatureConfig | None}, \code{dict[str, pd.DataFrame]}, \code{tuple[int, ...]} và tham
chiếu tới lớp chưa được định nghĩa xong, mà không tốn chi phí lúc chạy.\\
\code{from dataclasses import dataclass} & F, E & Không &
Decorator tự sinh \code{__init__}, \code{__repr__}, \code{__eq__} từ các trường khai báo. Dùng cho
\code{FeatureConfig} và \code{ResultsLog}.\\
\code{from dataclasses import field} & F, E & Có &
Dùng \code{field(default_factory=list)} để mỗi đối tượng có \emph{một danh sách riêng}. Viết
\code{rows: list = []} là lỗi trong dataclass, vì mọi đối tượng sẽ dùng chung một list có thể thay đổi
(\emph{mutable default}).\\
\code{from pathlib import Path} & F, E & Không &
Đường dẫn hướng đối tượng, ghép bằng toán tử \code{/}, chạy đúng trên cả Windows và Linux.\\
\code{import numpy as np} & F, E & Không &
Tính toán vector: \code{np.sin/cos} (mã hóa tuần hoàn), \code{np.abs}, \code{np.mean}, \code{np.sqrt},
\code{np.isfinite} (các chỉ số).\\
\code{import pandas as pd} & F, E & Không &
Xử lý bảng và chuỗi thời gian: đọc Parquet, \code{tz_convert}, \code{shift}, \code{rolling}, \code{groupby}.\\
\code{import matplotlib}; \code{matplotlib.use("Agg")} & E & Không &
Chọn backend \emph{Agg} (vẽ ra file, không mở cửa sổ) để script chạy được trên máy không có giao diện hoặc
server. Phải gọi \emph{trước} khi import \code{pyplot}.\\
\code{import matplotlib.pyplot as plt} & E & Không &
Vẽ các hình chẩn đoán (đường dự báo, phần dư, biểu đồ cột).\\
\bottomrule
\end{longtable}
\endgroup

\subsubsection{File \code{features.py}}\label{sec:impl-features}

\paragraph{Biến toàn cục.} Theo quy ước Python, tên viết HOA là hằng số; tên bắt đầu bằng dấu gạch dưới
(\code{_FIXED_HOLIDAYS}) là ``nội bộ'', không nên import từ bên ngoài.

\begingroup\footnotesize
\begin{longtable}{>{\raggedright\arraybackslash}p{0.25\linewidth}>{\raggedright\arraybackslash}p{0.27\linewidth}p{0.38\linewidth}}
\caption{Biến toàn cục trong \code{features.py}.}\label{tab:globals-features}\\
\toprule
Tên & Giá trị & Ý nghĩa\\
\midrule
\endfirsthead
\toprule
Tên & Giá trị & Ý nghĩa\\
\midrule
\endhead
\code{ROOT_DIR} & \code{Path(__file__).resolve()} \code{.parent.parent} &
Thư mục gốc repo, suy ra từ vị trí file. Nhờ vậy script chạy đúng dù được gọi từ thư mục nào.\\
\code{SPLIT_DIR} & \code{ROOT_DIR/"data"/"splits"} & Nơi phần A ghi \code{train/val/test.parquet}.\\
\code{TARGET} & \code{"total load actual"} & Tên cột mục tiêu $y_t$. Mọi mô-đun dùng hằng này thay vì gõ lại chuỗi, tránh lỗi chính tả.\\
\code{TSO_FORECAST} & \code{"total load forecast"} & Cột dự báo của TSO: giữ lại để làm mốc so sánh, \textbf{không} bao giờ là đặc trưng đầu vào.\\
\code{LOCAL_TZ} & \code{"Europe/Madrid"} & Múi giờ địa phương để tính đặc trưng lịch (CET/CEST).\\
\code{_FIXED_HOLIDAYS} & 9 chuỗi \code{"MM-DD"} & Các ngày lễ quốc gia có ngày cố định hằng năm.\\
\code{_GOOD_FRIDAYS} & 4 ngày 2015--2018 & Thứ Sáu Tuần Thánh (ngày di động theo lịch Phục Sinh).\\
\code{SPAIN_HOLIDAYS} & \code{pd.DatetimeIndex} (49 ngày) &
Ghép 9 ngày cố định $\times$ 5 năm (2014--2018) với 4 ngày Thứ Sáu Tuần Thánh. Dùng \code{.isin()} để gắn cờ ngày lễ rất nhanh.\\
\bottomrule
\end{longtable}
\endgroup

\paragraph{Lớp \code{FeatureConfig}} (\code{@dataclass}) gom toàn bộ ``siêu tham số'' của pipeline đặc trưng
vào một đối tượng, để thí nghiệm ablation và so sánh chân trời chỉ cần đổi cấu hình, không phải sửa code.

\begingroup\footnotesize
\begin{longtable}{>{\raggedright\arraybackslash}p{0.22\linewidth}>{\raggedright\arraybackslash}p{0.25\linewidth}p{0.43\linewidth}}
\caption{Các trường của \code{FeatureConfig}.}\label{tab:featureconfig}\\
\toprule
Trường & Mặc định & Ý nghĩa\\
\midrule
\endfirsthead
\toprule
Trường & Mặc định & Ý nghĩa\\
\midrule
\endhead
\code{horizon} & \code{24} & Chân trời dự báo $h$ (giờ). 24 là dự báo ngày tới, 1 là dự báo giờ tới.\\
\code{lags} & \code{(1,2,3,24,48,72,168,336)} & Danh sách độ trễ \emph{ứng viên}; chỉ những lag $\geq h$ được dùng thật.\\
\code{roll_windows} & \code{(24, 168)} & Độ dài cửa sổ trượt (1 ngày, 1 tuần) cho trung bình và độ lệch chuẩn.\\
\code{comfort_temp_c} & \code{18.0} & Nhiệt độ gốc (°C) cho độ-giờ sưởi; ngưỡng làm mát là \code{comfort_temp_c + 4} = 22\,°C.\\
\code{use_calendar}, \code{use_lags}, \code{use_weather} & \code{True} & Bật/tắt từng nhóm đặc trưng, dùng cho ablation (\cref{tab:ablation}).\\
\code{extra_drop} & \code{[]} (qua \code{field}) & Tên các cột muốn bỏ thêm khi thử nghiệm.\\
\bottomrule
\end{longtable}
\endgroup

Thuộc tính \code{effective_lags} (\code{@property}, tính lại mỗi lần truy cập) trả về
\texttt{sorted(\{lag for lag in lags if lag >= horizon\} | \{horizon\})}. Đây là chốt chặn rò rỉ quan trọng nhất của
phần B: với $h=24$ kết quả là $[24, 48, 72, 168, 336]$ (lag 1--3 tự động bị loại); với $h=1$ là
$[1, 2, 3, 24, 48, 72, 168, 336]$. Phép hợp \texttt{| \{horizon\}} bảo đảm luôn có cột \code{lag_h}, vì mô hình
cơ sở ngây thơ $\hat y_t = y_{t-h}$ cần cột này.

\paragraph{Các hàm.}

\begingroup\footnotesize
\begin{longtable}{>{\raggedright\arraybackslash}p{0.31\linewidth}p{0.63\linewidth}}
\caption{Các hàm trong \code{features.py}.}\label{tab:functions-features}\\
\toprule
Hàm & Chức năng và cách hoạt động\\
\midrule
\endfirsthead
\toprule
Hàm & Chức năng và cách hoạt động\\
\midrule
\endhead
\code{load_splits(split_dir)} &
\textbf{Vào:} thư mục chứa các tập đã chia. \textbf{Ra:} \code{dict} gồm 3 DataFrame \code{"train"}, \code{"val"},
\code{"test"}. Đọc file Parquet nên giữ nguyên kiểu \code{datetime64[UTC]} của cột \code{time}.\\
\addlinespace
\code{add_calendar_features(df)} &
\textbf{Vào:} DataFrame có cột \code{time} (UTC). \textbf{Ra:} DataFrame 15 cột \code{cal_*}, cùng index.
Các bước: (1) \code{tz_convert(LOCAL_TZ)} đổi sang giờ địa phương; (2) \code{tz_localize(None).dt.normalize()}
lấy ngày địa phương (00:00) để so với \code{SPAIN_HOLIDAYS}; (3) trích giờ/thứ/tháng/ngày trong năm;
(4) tạo cờ cuối tuần, ngày lễ, ngày nghỉ; (5) ngày ``cầu'' = ngày làm việc mà hôm trước \emph{và} hôm sau
đều là ngày nghỉ, kiểm tra bằng cách dịch ngày $\pm 1$ và dùng \texttt{(dow $\pm$ 1) \% 7}; (6) giờ mùa hè được
suy ra từ độ lệch giờ địa phương so với UTC (2 giờ = CEST); (7) mã hóa $\sin/\cos$ cho giờ, thứ, ngày trong năm.\\
\addlinespace
\code{add_lag_features(df, cfg)} &
\textbf{Vào:} DataFrame có cột \code{TARGET} đã sắp theo thời gian. \textbf{Ra:} 13 cột (khi $h=24$).
\code{y.shift(k)} dịch chuỗi xuống $k$ dòng, nên dòng $t$ nhận $y_{t-k}$. Biến \code{base = y.shift(h)}
là giá trị cuối cùng quan sát được, mọi cửa sổ \code{rolling} đều tính trên \code{base}, nên không bao giờ chứa
$y_t$. \code{min_periods=w} bắt buộc cửa sổ đủ dài, nếu không thì trả NaN (sau đó bị loại).
\code{lag_same_hour_4w_mean} ghép 4 chuỗi dịch $168k$ giờ bằng \code{pd.concat(axis=1)} rồi lấy trung bình
theo hàng với \code{skipna=False} (thiếu một tuần thì kết quả là NaN). \code{diff_24} $= y_{t-h} - y_{t-h-24}$.\\
\addlinespace
\code{add_weather_features(df, cfg)} &
\textbf{Vào:} các cột thời tiết do phần A tổng hợp. \textbf{Ra:} 11 cột \code{w_*}. Đổi Kelvin sang °C
($-273.15$); \code{clip(lower=0)} hiện thực hàm $\max(0,\cdot)$ của HDH/CDH; \code{rolling(24, min_periods=1)}
cho nhiệt độ trung bình 24 giờ qua (\code{min_periods=1} để không sinh NaN ở đầu chuỗi).\\
\addlinespace
\code{build_feature_frame(full, cfg)} &
\textbf{Vào:} toàn bộ dữ liệu (train+val+test nối lại). \textbf{Ra:} một DataFrame gồm \code{time},
\code{TARGET}, \code{TSO_FORECAST} và các đặc trưng. Sắp xếp theo thời gian, \code{assert} thời gian không trùng,
gọi ba hàm \code{add_*} theo các cờ \code{use_*}, ghép theo cột bằng \code{pd.concat(axis=1)}, rồi bỏ các cột
trong \code{extra_drop}.\\
\addlinespace
\code{feature_columns(frame)} &
\textbf{Ra:} danh sách tên cột đầu vào cho mô hình, tức mọi cột trừ \code{time}, \code{TARGET},
\code{TSO_FORECAST}. Đây là \emph{nguồn duy nhất} quyết định mô hình được thấy cột nào, nhờ đó các phần C và D
dùng cùng một bộ đặc trưng.\\
\addlinespace
\code{make_datasets(cfg, split_dir)} &
\textbf{Hàm chính, dùng chung cho cả nhóm.} \textbf{Ra:} \code{dict} gồm 3 DataFrame đã có đặc trưng.
(1) Ghi lại mốc thời gian đầu/cuối của từng tập; (2) nối 3 tập theo thời gian và gọi
\code{build_feature_frame}, để lag của những giờ đầu tập val lấy được giá trị cuối tập train (là quá khứ hợp lệ);
(3) cắt lại theo các mốc thời gian; (4) \code{dropna} trên các cột đặc trưng, tức bỏ 672 giờ khởi động đầu tập
train; (5) hai \code{assert} kiểm tra các tập vẫn tách rời theo thời gian và val/test \emph{không mất dòng nào}.\\
\addlinespace
\code{xy(df, cols)} &
Hàm tiện ích tách DataFrame thành \code{(X, y)}. \code{train_models.py} hiện truy cập cột trực tiếp, nên hàm
này dành cho phần D/E nếu cần.\\
\addlinespace
Khối \code{if __name__ == "__main__"} &
Chỉ chạy khi gọi \code{python src/features.py}, không chạy khi file được import. In ra chân trời dự báo,
danh sách 39 đặc trưng và kích thước từng tập, dùng để kiểm tra nhanh.\\
\bottomrule
\end{longtable}
\endgroup

\subsubsection{File \code{evaluate.py}}\label{sec:impl-evaluate}

\paragraph{Biến toàn cục.} Chỉ có một hằng số: \code{METRICS = ("MAE", "RMSE", "MAPE", "sMAPE", "R2", "MaxAE")}.
Đây là danh sách tên chỉ số theo đúng thứ tự xuất trong bảng. Các mô-đun khác import hằng này khi cần lặp qua
mọi chỉ số (ví dụ để vẽ hoặc định dạng bảng), thay vì tự gõ lại. Dùng \code{tuple} thay vì \code{list} để
không thể bị sửa nhầm lúc chạy.

\paragraph{Các hàm tính chỉ số.} Tất cả có chung chữ ký \code{f(y_true, y_pred) -> float}, nhận bất kỳ dạng
mảng nào (list, \code{np.ndarray}, \code{pd.Series}) và trả về một số \code{float} của Python (không phải
\code{np.float64}), để ghi ra CSV/JSON gọn gàng.

\begingroup\footnotesize
\begin{longtable}{>{\raggedright\arraybackslash}p{0.31\linewidth}p{0.63\linewidth}}
\caption{Các hàm và lớp trong \code{evaluate.py}.}\label{tab:functions-evaluate}\\
\toprule
Hàm / lớp & Chức năng và cách hoạt động\\
\midrule
\endfirsthead
\toprule
Hàm / lớp & Chức năng và cách hoạt động\\
\midrule
\endhead
\code{_as_arrays(y_true, y_pred)} &
Hàm nội bộ, được \emph{mọi} hàm chỉ số gọi đầu tiên. Chuyển đầu vào sang \code{np.float64} và làm phẳng
(\code{ravel}), sau đó \code{raise ValueError} nếu hai mảng khác kích thước hoặc chứa NaN/inf. Nhờ vậy một lỗi
như quên \code{dropna} sẽ báo ngay, thay vì âm thầm cho ra MAE bằng NaN hoặc sai lệch.\\
\addlinespace
\code{mae}, \code{rmse} & $\frac1n\sum|e_i|$ và $\sqrt{\frac1n\sum e_i^2}$ với $e_i = y_i - \hat y_i$.\\
\addlinespace
\code{mape} & $\frac{100}{n}\sum|e_i/y_i|$. Báo lỗi nếu có $y_i = 0$ (phép chia không xác định).\\
\addlinespace
\code{smape} & MAPE đối xứng, mẫu số là $(|y_i| + |\hat y_i|)/2$, nên dự báo cao hay thấp đều bị phạt như nhau.\\
\addlinespace
\code{r2} & $1 - SS_{res}/SS_{tot}$. Bằng 1 khi dự báo hoàn hảo, bằng 0 khi chỉ đoán bằng giá trị trung bình,
và có thể âm. Lưu ý: không xác định khi $y$ là hằng số ($SS_{tot}=0$); điều này không xảy ra với phụ tải điện.\\
\addlinespace
\code{max_ae} & $\max_i |e_i|$: sai số của giờ tệ nhất.\\
\addlinespace
\code{evaluate(y_true, y_pred)} &
Gọi cả 6 hàm và trả về \code{dict} \texttt{\{"MAE": ..., ..., "MaxAE": ...\}}. Đây là điểm vào duy nhất cần nhớ.\\
\addlinespace
Lớp \code{ResultsLog} &
\code{@dataclass} có một trường \code{rows: list[dict]} (khởi tạo bằng \code{field(default_factory=list)}),
mỗi phần tử là một dòng kết quả. Các phương thức:
\begin{itemize}[leftmargin=1em,nosep]
  \item \code{add(model, split, y_true, y_pred, **extra)}: gọi \code{evaluate}, ghép thêm tên mô hình, tên tập và
        các thông tin phụ tùy ý qua \code{**extra} (ví dụ \code{alpha}, \code{fit_s}, \code{family}), lưu vào
        \code{rows} và trả về \code{dict} chỉ số.
  \item \code{to_frame()}: chuyển \code{rows} thành \code{pd.DataFrame}.
  \item \code{pivot(metric)}: bảng mô hình $\times$ tập cho một chỉ số, ví dụ MAE theo val/test.
  \item \code{save(path)}: tạo thư mục nếu chưa có, ghi CSV với 4 chữ số thập phân.
\end{itemize}\\
\addlinespace
\code{error_by_group(time, y_true, y_pred, by, tz)} &
Tính MAE theo nhóm \code{"hour"}, \code{"dayofweek"} hoặc \code{"month"} theo \emph{giờ địa phương}. Đổi múi
giờ, chọn khóa nhóm bằng một \code{dict} tra cứu, tạo \code{Series} sai số tuyệt đối có index là khóa, rồi
\code{groupby(level=0).mean()}. Dùng để vẽ \cref{fig:mae-hour}.\\
\addlinespace
\code{plot_forecast(...)} &
Vẽ đường thực tế (đen, đậm) và dự báo của nhiều mô hình trong một cửa sổ \code{days} ngày bắt đầu từ
\code{start}, lọc bằng mặt nạ boolean trên cột thời gian. Dùng cho \cref{fig:forecast}.\\
\addlinespace
\code{plot_residuals(...)} &
Hình 2 ô: histogram phần dư $y - \hat y$ (tiêu đề ghi trung bình và độ lệch chuẩn), và biểu đồ phân tán dự báo
so với thực tế, kèm đường chéo $y=\hat y$ màu đỏ. Dùng cho \cref{fig:residuals}.\\
\addlinespace
\code{plot_metric_bars(table, metric, out_path, split)} &
Biểu đồ cột ngang một chỉ số cho mọi mô hình trên một tập. Sắp tăng dần (riêng $R^2$ sắp giảm dần vì càng lớn
càng tốt), ghi giá trị cạnh mỗi cột. Dùng cho \cref{fig:bars}.\\
\bottomrule
\end{longtable}
\endgroup

Mọi hàm vẽ đều tự tạo thư mục đích (\code{mkdir(parents=True, exist_ok=True)}), lưu ở 150 dpi và gọi
\code{plt.close(fig)} để giải phóng bộ nhớ khi vẽ nhiều hình liên tiếp.

\subsubsection{Các quyết định thiết kế chung}\label{sec:impl-design}
\begin{itemize}
  \item \textbf{Hàm thuần (pure function):} các hàm \code{add_*} không sửa DataFrame đầu vào mà trả về DataFrame
        mới có cùng index. Nhờ vậy có thể ghép bằng \code{pd.concat(axis=1)} và test riêng từng hàm.
  \item \textbf{Cấu hình tách khỏi code:} mọi tham số nằm trong \code{FeatureConfig}. Thí nghiệm ablation chỉ
        là một vòng lặp qua các cấu hình khác nhau.
  \item \textbf{Báo lỗi sớm:} \code{assert} trong \code{features.py} kiểm tra các bất biến của pipeline (không
        trùng thời gian, không chồng lấn tập, không mất dòng val/test). \code{ValueError} trong
        \code{evaluate.py} kiểm tra dữ liệu do người dùng truyền vào, và vẫn hoạt động khi chạy Python với
        cờ \code{-O}.
  \item \textbf{Một nguồn sự thật:} tên cột (\code{TARGET}, \code{TSO_FORECAST}), danh sách đặc trưng
        (\code{feature_columns}) và cách tính chỉ số (\code{evaluate}) chỉ định nghĩa một lần, nên kết quả của
        các phần C, D, E so sánh trực tiếp được với nhau.
\end{itemize}
